\documentclass[11pt,addpoints]{exam}   	% use "amsart" instead of "article" for AMSLaTeX format
\usepackage{geometry}                		% See geometry.pdf to learn the layout options. There are lots.
\geometry{letterpaper}                   		% ... or a4paper or a5paper or ... 
%\geometry{landscape}                		% Activate for rotated page geometry
%\usepackage[parfill]{parskip}    		% Activate to begin paragraphs with an empty line rather than an indent
\usepackage{graphicx}				% Use pdf, png, jpg, or eps§ with pdflatex; use eps in DVI mode
								% TeX will automatically convert eps --> pdf in pdflatex		
\usepackage{amsmath}
\usepackage{cancel}
\usepackage{amssymb}
\usepackage{multicol}
\usepackage{mhchem}
\usepackage{lewis}

\title{Quarter 3 Exam Review}
\author{Mark Peever}
\date{March 6, 2026}							% Activate to display a given date or no date

\begin{document}
\maketitle

\pointsinrightmargin
%\marginpointname{ \points}
\printanswers

\begin{center}
\fbox{\fbox{\parbox{5.5in}{\centering
Answer all questions on a separate sheet of paper. Show all your work.
}}}
\end{center}
\vspace{0.1in}
\makebox[\textwidth]{Name:\enspace\hrulefill}
\vspace{0.2in}

\begin{questions}

\question[5] Give the electron configurations for:
\begin{itemize}
\item \ce{He}
\item \ce{O}
\item \ce{S}
\item \ce{Cl}
\end{itemize}

\begin{solution}
\begin{itemize}
\item $1s^{2} $
\item $1s^{2} 2s^{2} 2p^{4} $
\item $1s^{2} 2s^{2} 2p^{6} 3s^{2} 3p^{4}$
\item $1s^{2} 2s^{2} 2p^{6} 3s^{2} 3p^{5}$
\end{itemize}
\end{solution}

\clearpage
\question[5] 
Give the Lewis Structures for:
\begin{itemize}
\item \ce{H}
\item \ce{HCl}
\item \ce{H2O}
%\item \ce{NaOH}
\end{itemize}

\begin{solution}
\begin{itemize}
\item \lewis{H}{}{}{}{}{}{.}{}{}
\item \lewis{H}{}{}{}{}{}{}{}{}\hspace{-0.5em}--\hspace{-0.75em}\lewis{Cl}{}{}{.}{.}{.}{.}{.}{.}
\item \lewis{H}{}{}{}{}{}{}{}{}\hspace{-0.5em} -- \hspace{-0.75em}\lewis{O}{}{}{.}{.}{}{}{.}{.}\hspace{-0.5em} -- \hspace{-0.75em}\lewis{H}{}{}{}{}{}{}{}{}
%\item \lewis{\ce{Na^{+}}}{}{}{}{}{}{}{}{} \lewis{ O}{.}{.}{.}{.}{}{}{.}{.} \hspace{-0.5em}--\hspace{-0.75em}\lewis{H}{}{}{}{}{}{}{}{}
\end{itemize}
\end{solution}

\question[5] What is the correct chemical equation to show \ce{HCl} producing \ce{H3O^{+}} ions in water?

\begin{solution}
\ce{HCl (aq) + H2O (l) -> H3O^{+} (aq) + Cl^{-} (aq)}
\end{solution}

\clearpage
\question[5] How many \emph{moles} of \ce{HCl} are in a $124.00 mL$ sample of $0.250 M$ solution?

\begin{solution}
\begin{equation} 
\begin{split}
    [HCl] &= 0.250 M \\
    V_{HCl} &= 0.12400 L\\
    n &= ?\\
    [HCl] = \frac{n_{HCl}}{V_{HCl}} \\
    n_{HCl} &= [HCl] \cdot V_{HCl} \\
                 &= (0.250 \frac{mol}{L})(0.12400 L) \\
                 &= 0.0310 mol \\
 \end{split}
 \end{equation}
\end{solution}

\question[5]  What is the molar mass of \ce{HCl}?

\begin{solution}
\begin{equation} 
\begin{split}
    M_{HCl} &= M_{H} + M_{Cl} \\
                  &= (1.00794 \frac{g}{mol}) + (35.453 \frac{g}{mol}) \\
                  &= 35.46094 \frac{g}{mol} \\
                  &= 35.461 \frac{g}{mol} \\
 \end{split}
 \end{equation}
\end{solution}

\question[5]  How many grams of \ce{HCl} are in $0.100 L$ of $2.34 M$ solution?

\begin{solution}
\begin{equation} 
\begin{split}
    m &= n_{HCl} \cdot M_{HCl} \\
        &= ([HCl])(V_{HCl}) \cdot M_{HCl} \\
        &= (0.100 L) (2.34 \frac{mol}{L}) (35.461 \frac{g}{mol}) \\
        &= 8.297874 mol \\
        &= 8.30 mol
 \end{split}
 \end{equation}
\end{solution}

\question[5] What is the volume of $0.1000 M$ solution you get from diluting $10.00 mL$ of $3.500 M$ solution?

\begin{solution}
\begin{equation} 
\begin{split}
   V_1 [ \hspace{0.1 in}]_1 &= V_2 [ \hspace{0.1 in}]_2 \\
   V_2 &= \frac{V_1 [ \hspace{0.1 in}]_1} {[\hspace{0.1 in} ]_2} \\
          &= \frac{(10.00 mL) (3.500 \frac{mol}{L}) } {0.1000 \frac{mol}{L}} \\
          &= 350.0 mL
 \end{split}
 \end{equation}
\end{solution}



\end{questions}

\pagebreak
\section*{Fun Facts}
\begin{multicols}{2}
\begin{itemize}
\item 
$ [\hspace{0.2 in}] = \frac{n}{V}$

\item 
$ [\hspace{0.2 in}]_1 \cdot V_1=  [\hspace{0.2 in}]_2 \cdot V_2$

\item
$q = m c \Delta T$
\vspace{0.2in}

\item
$\Delta T = T_{final} - T_{initial}$
\vspace{0.2in}

\item
$\rho = \frac{m}{V} $
\vspace{0.2in}

\item
$T_K = T_C + 273.15 K$
\vspace{0.2in}

\item
$T_F = \frac{9}{5}T_C + 32^{\circ}F$
\vspace{0.2in}

\item
$T_C = \frac{5}{9}(T_F - 32^{\circ}F)$
\vspace{0.2in}

\item
$c_{water} = 1.000 \frac{calorie}{g \cdot ^{\circ}C}$
\vspace{0.2in}

\item
$c_{water} = 4.184 \frac{J}{g \cdot ^{\circ}C}$
\vspace{0.2in}

\item
$c_{copper} = 0.3851 \frac{J}{g \cdot ^{\circ}C}$
\vspace{0.2in}

\item
$\rho_{water} = 1.000 \frac{g}{cm^3}$
\vspace{0.2in}

\item
$\rho_{copper} = 8.96 \frac{g}{cm^3}$
\vspace{0.2in}

\item
$ 1.0000 cm^3 = 1.0000 mL $
\vspace{0.2in}

\item
$ 1.000 cal = 4.184 J $
\vspace{0.2in}

\item
$m = n \times M$
\vspace{0.2in}

\item
$m = n \times m_{molar}$
\vspace{0.2in}

\item
$1 amu = 1 \frac{g}{mol}$ 
\vspace{0.2in}

\item
$N_A = 6.02214076 \times 10^{23} $ 
\vspace{0.2in}

\end{itemize}
\vspace{0.2in}

\end{multicols}




\pagebreak

\begin{table}
\centering
\begin{tabular}[p]{| l | l |}
\textbf{Prefix} & \textbf{Count} \\
\hline
mono & one \\
di       & two \\
tri       & three \\
tetra       & four \\
penta       & five \\
hecta       & six \\
hepta       & seven \\
octa       & eight \\
nona       & nine \\
deca       & ten \\
\end{tabular}
\caption{Prefixes for Covalent Compound Names}
\label{table:covalent-prefix}
\end{table}






\end{document}  